\documentclass[11pt]{article}
\usepackage[letterpaper,margin=1in]{geometry}
\usepackage{amsmath,amssymb,amsthm,mathtools}
\usepackage{booktabs,array}
\usepackage{microtype}
\usepackage[hidelinks]{hyperref}
\usepackage{xcolor}
\hypersetup{
  pdftitle={The Riemann Kernel Is Globally PF4},
  pdfauthor={Joshuah Rainstar},
  pdfsubject={Strict global PF4 classification of the Riemann kernel},
  pdfkeywords={Polya-frequency function, total positivity, Riemann kernel, entire function, Wronskian, Fourier transform, exact certificate}
}

\newtheorem{theorem}{Theorem}[section]
\newtheorem{proposition}[theorem]{Proposition}
\newtheorem{lemma}[theorem]{Lemma}
\newtheorem{corollary}[theorem]{Corollary}
\theoremstyle{remark}
\newtheorem{remark}[theorem]{Remark}

\newcommand{\R}{\mathbb{R}}
\newcommand{\T}{\mathcal{T}}
\newcommand{\Cfour}{\mathcal{C}_4}
\newcommand{\E}{\mathbb{E}}
\newcommand{\dd}{\,\mathrm{d}}
\newcommand{\PF}{\mathrm{PF}}
\newcommand{\TP}{\mathrm{TP}}
\newcommand{\TN}{\mathrm{TN}}
\newcommand{\secref}[1]{Section~\ref{#1}}
\newcommand{\appref}[1]{Appendix~\ref{#1}}

\numberwithin{equation}{section}

\title{The Riemann Kernel Is Globally \texorpdfstring{$\PF_4$}{PF4}}
\author{Joshuah Rainstar\thanks{Email:
\href{mailto:joshuah.rainstar@gmail.com}{\texttt{joshuah.rainstar@gmail.com}}}}
\date{Revised 22 July 2026}

\begin{document}
\maketitle
\begin{abstract}
Let \(\Phi\) be the classical Riemann, or de Bruijn--Newman, kernel. We prove
that its exact global P\'olya-frequency order is four: every translation minor
of order at most four is strictly positive at strictly ordered nodes, whereas
the order-five minor at the exact spacing \(211/2000\) is negative. For the
positive half, a two-mode
argument proves global positivity of the log-curvature quantities \(q,F_2\)
and of the fully confluent invariant \(\Cfour\). The finite verification uses
only rational polynomial algebra, Taylor inequalities, and geometric tails.
The direct finite-witness certificate likewise uses exact rational arithmetic.
A
quotient-Wronskian argument reduces arbitrary fourth-order minors to
\(\partial_\xi\Psi<0\). In the increasing curvature coordinate, its numerator
is
\[
 \delta\Lambda\int_p^w
 \Delta(t)\frac{\Cfour(t)}{Q(t)^6\kappa(t)^2}\dd t.
\]
Here \(\Delta\) is a deterministic closed cumulative gap built from two
explicitly normalized, curvature-tilted triangular weights. Its strict
positivity follows on the coordinate interval from the exact mass identities
and one explicit algebraic sign change. The transport
reductions and analytic sign estimates are given explicitly, with a table of
the finite polynomial inequalities used for the three global signs.
\end{abstract}
\noindent\textbf{2020 Mathematics Subject Classification.}
15B48 (primary); 42A38, 30D15, 11M26 (secondary).

\noindent\textbf{Keywords.}
P\'olya-frequency function, total positivity, Riemann kernel, entire function,
Wronskian, Fourier transform, exact certificate.
\tableofcontents

\section{Introduction}\label{sec:introduction}

A function \(f:\R\to\R\) is \(\PF_r\) when every translation-kernel minor
\[
 \det[f(x_i-y_j)]_{i,j=1}^k
\]
is nonnegative for \(1\leq k\leq r\) and increasing node tuples. We call the
property strict when the determinant is positive whenever both tuples are
strictly increasing. This paper determines the exact finite P\'olya-frequency
order of the Riemann kernel; see \cite{belton} for modern background and
Schoenberg's characterization. Related finite-order and Riemann-kernel work is
placed in \secref{sec:related-work}.

\begin{theorem}[Exact global PF order]\label{thm:main}
The exact global P\'olya-frequency order of \(\Phi\) is four. More precisely:
\begin{enumerate}
\item for every \(1\leq k\leq4\) and strictly increasing real tuples
\(x_1<\cdots<x_k\) and \(y_1<\cdots<y_k\),
\[
 \det[\Phi(x_i-y_j)]_{i,j=1}^k>0.
\]
\item with \(h=211/2000\),
\[
 \det[\Phi((i-j)h)]_{i,j=0}^4<0.
\]
\end{enumerate}
\end{theorem}

The two theorem inputs are deliberately separated. The positive input is the
universal strict-order theorem through four. The negative input is the direct
rational order-five witness in part~(2), including proved strict ordering of
its nodes and the signed Toeplitz-to-translation-matrix identification.

The proof has two one-variable inputs and one exact transport
mechanism. First, global inequalities \(q>0\) and \(F_2>0\) imply strict
\(\PF_3\) and positivity of a three-point functional \(\Lambda\). Second, the
fully confluent order-four invariant \(\Cfour\) is globally positive. Exact
quotient identities reduce strict \(\PF_4\) to the monotonicity of a scalar
three-point function \(\Psi\). On the actual image of an increasing curvature
coordinate, the derivative of \(\Psi\) becomes a positive integral of
\(\Cfour\) against a deterministic closed cumulative gap.

The finite part of the positive input consists of exact rational Bernstein
coefficients on fixed algebraic domains, followed by analytic bounds for all
remaining theta modes; \appref{app:tail} lists the transformed target
polynomials, domains, degrees, minimum coefficients, derivative envelopes,
and the surviving relative margins. The transport algebra is proved in the
text, and the order-five
calculation uses rational enclosures and elementary series bounds.
\section{Relation to earlier work}\label{sec:related-work}

Schoenberg introduced P\'olya-frequency functions and connected translation
minors with bilateral Laplace transforms and variation-diminishing convolution
operators~\cite{schoenberg-ii,belton}. Karlin's monograph remains the
standard systematic account of total positivity and its composition
principles~\cite{karlin}; a modern review of PF functions and their preservers
is given by Belton, Guillot, Khare, and Putinar~\cite{belton}.

Finite-order translation kernels have a substantially wider range than the
infinite-order class. Khare's characterization of measurable \(\TN_p\)
functions for \(p\geq3\) provides a general finite-order setting~\cite{khare}.
The present proof instead derives a criterion specific to order four, using
successive quotient derivatives and a curvature transport identity.

The theta normalization used here is the classical one in de Bruijn's study of
the Riemann Fourier kernel~\cite{debruijn}. For the same kernel, Dimitrov and Xu
relate Wronskians of Fourier and Laplace transforms to correlation kernels and
real-zero criteria~\cite{dimitrov-xu}, while Csordas and Varga obtain strict
concavity and moment inequalities connected with the Riemann
Hypothesis~\cite{csordas-varga}. Grochenig explains the distinct Schoenberg
correspondence relevant to the Riemann Hypothesis: the PF function there is
recovered from the reciprocal of a Laguerre--P\'olya entire function, rather than
identified with the original Riemann kernel~\cite{grochenig}.

Micha\l{}owski first published an explicit negative \(5\times5\) Toeplitz
minor for this kernel, supported by an interval computation, together with
single-configuration positivity checks at orders at most
four~\cite{michalowski}. Global \(\PF_4\) is posed there as Open Problem 1.
Theorem~\ref{thm:main} resolves that problem and supplies an exact rational
negative order-five minor, completing the finite-order classification.
\section{The kernel and smooth even continuation}\label{sec:kernel}

Put
\[
 \vartheta(x)=\sum_{n\in\mathbb Z}e^{-\pi n^2x},\qquad
 \Theta(t)=\vartheta(e^{2t}),\qquad H(t)=e^{t/2}\Theta(t).
\]
This is a real-valued definition: summability is first proved directly for
every \(x>0\), before any analytic continuation is invoked. The real
specialization of Gaussian Poisson summation gives the exact normalization
\(\vartheta(x)=x^{-1/2}\vartheta(1/x)\), and hence
\[
 \Theta(-t)=e^t\Theta(t),\qquad H(-t)=H(t).
\]
A Jacobi-theta realization supplies a globally analytic function agreeing
with this real series. Thus \(H\) is real analytic and even. Define
\[
 \Phi(t)=\frac12\left(H''(t)-\frac14H(t)\right).
\]
This is twice the common de Bruijn kernel normalization; multiplication by a
positive constant changes none of the PF signs or strictness statements.
On \(t\geq0\), locally uniform polynomial--geometric majorants justify
termwise differentiation through order six and give
\begin{equation}\label{eq:kernel-series}
 \Phi(t)=\sum_{n\geq1}
 \left(4\pi^2n^4e^{9t/2}-6\pi n^2e^{5t/2}\right)
 e^{-\pi n^2e^{2t}}.
\end{equation}
Parity reflects this identity to every real \(t\), equivalently by replacing
\(t\) on the right by \(|t|\). Thus the positive-side formula is the
restriction of a real-analytic even function. Smoothness and every derivative
used later at the origin are consequences of that analytic realization, not
of numerical coverage or of differentiating an absolute-value definition.

For \(t\geq0\), the coefficient of the \(n\)-th exponential can be factored as
\[
 2\pi n^2e^{5t/2}\bigl(2\pi n^2e^{2t}-3\bigr)>0,
\]
because \(2\pi>3\). Hence every term in \eqref{eq:kernel-series} is positive
and \(\Phi(t)>0\) on \([0,\infty)\). Evenness gives \(\Phi(t)>0\) on all of
\(\R\), so the following logarithm is globally defined.

Write
\[
 \ell=\log\Phi,\qquad s=\ell',\qquad q=-s'=-\ell''.
\]
For \(a<b\), define
\begin{equation}\label{eq:AM}
 A(a,b)=s(a)-s(b),\qquad
 M(a,b)=\frac{q(b)-q(a)}{A(a,b)}.
\end{equation}
When \(q>0\), \(s\) is strictly decreasing and \(A(a,b)>0\).
\section{Certified one-variable inputs}\label{sec:inputs}

Define
\[
 F_1=qq''-(q')^2,\qquad F_2=q^3-F_1.
\]
For the logarithmic cumulants \(c_j=\ell^{(j)}\), set
\[
\begin{aligned}
 &m_0=1,\quad m_1=0,\quad m_2=c_2,\quad m_3=c_3,\\
 &m_4=c_4+3c_2^2,\quad m_5=c_5+10c_2c_3,\\
 &m_6=c_6+15c_2c_4+10c_3^2+15c_2^3,
\end{aligned}
\]
and
\begin{equation}\label{eq:C4-def}
 \Cfour(t)=\det[m_{i+j-2}(t)]_{i,j=1}^4.
\end{equation}

\begin{lemma}[Certified input]\label{lem:certified-input}
For the kernel \(\Phi\),
\[
 q(t)>0,\qquad F_2(t)>0,\qquad \Cfour(t)>0
 \quad(t\in\R).
\]
\end{lemma}

The proof combines exact rational polynomial inequalities with analytic
theta-tail bounds. For \(u\geq0\), put
\[
 x=\pi e^{2u},\qquad \eta=e^{-3x},\qquad R_0(z)=2z-3,
\]
and recursively
\[
 R_{j+1}(z)=\left(\frac52-2z\right)R_j(z)+2zR_j'(z).
\]
After division by the positive first theta mode, the first two modes give the
normalized jet entries
\[
 a_j(x,\eta)=\frac{R_j(x)+4\eta R_j(4x)}{R_0(x)}.
\]
Exact rational Bernstein coefficients on three fixed algebraic domains prove
\[
 q_{1,2}>10,\qquad (F_2)_{1,2}>1000,
 \qquad (\Cfour)_{1,2}>50{,}000{,}000.
\]
Here the half-line certificates are applied directly to the cleared
numerators of the shifted targets \(q_{1,2}-10\),
\((F_2)_{1,2}-1000\), and
\((\Cfour)_{1,2}-50{,}000{,}000\). The complete contribution of modes
\(n\geq3\) is bounded by one monotone polynomial-geometric estimate. Its
changes in the same cleared sign numerators are smaller than the three
displayed margins, both on
\(\pi\leq x\leq5\) and on \(x\geq5\). Evenness completes the real line.
All inequalities, including the derivative envelopes used on the tail, are
derived in \appref{app:tail}. Exact raw-jet identities remove the positive
denominators, and the reflected global jet from \secref{sec:kernel} covers the
negative half-line and the origin.
\section{Order three and the positive slope functional}\label{sec:pf3}

For \(\xi<m<r\), define
\begin{equation}\label{eq:Lambda}
 \Lambda(\xi;m,r)=A(\xi,r)+M(\xi,m)-M(m,r).
\end{equation}
We give the full estimate used to prove its positivity.

Set \(f_0=q'/q\). Because \(s'=-q\),
\[
 M(a,b)=\frac{\int_a^bq'(t)\dd t}{\int_a^bq(t)\dd t}
       =\frac{\int_a^b f_0(t)q(t)\dd t}{\int_a^bq(t)\dd t};
\]
thus \(M(a,b)\) is a \(q\)-weighted mean of \(f_0\). Choose a point
\(u\in[\xi,m]\) where \(f_0\) is minimal on the left interval and a point
\(v\in[m,r]\) where it is maximal on the right. The weighted-mean bounds and
the ordering of the extremizers are
\[
 M(\xi,m)\geq f_0(u),\qquad M(m,r)\leq f_0(v),
 \qquad u\leq m\leq v.
\]
Consequently
\[
 M(m,r)-M(\xi,m)\leq f_0(v)-f_0(u)
 \leq\int_\xi^r\max(f_0'(t),0)\dd t.
\]
Since
\[
 f_0'=\frac{qq''-(q')^2}{q^2}=\frac{F_1}{q^2}
\]
and \(A(\xi,r)=\int_\xi^r q(t)\dd t\), equation
\eqref{eq:Lambda} yields
\begin{align}
 \Lambda(\xi;m,r)
 &\geq\int_\xi^r\left(q-\frac{\max(F_1,0)}{q^2}\right)\dd t\\
 &=\int_\xi^r\frac{q^3-\max(F_1,0)}{q^2}\dd t
 =\int_\xi^r\frac{\min(q^3,F_2)}{q^2}\dd t>0.
 \label{eq:Lambda-lower}
\end{align}
The second equality uses
\(q^3-\max(F_1,0)=\min(q^3,q^3-F_1)=\min(q^3,F_2)\); the last inequality
follows from Lemma~\ref{lem:certified-input} and \(q>0\).

The same quantity is the logarithmic slope derivative in the normalized
order-three determinant. Hence \eqref{eq:Lambda-lower}, together with strict
log-concavity, proves strict global \(\PF_3\). We rederive the needed
Wronskian form in the next section, where the quotient integral transfers
these confluent signs to every ordered collocation minor.
\section{Quotient-Wronskian reduction}\label{sec:quotient}

Fix \(y_1<y_2<y_3<y_4\), put \(u_j(t)=\Phi(t-y_j)\), and write
\(p_j=t-y_j\), so \(p_4<p_3<p_2<p_1\). Define
\[
 v_j=(u_j/u_1)',\qquad w_j=(v_j/v_2)'.
\]
These quotients are introduced only after their denominators are known to be
positive: \(u_1>0\) by \secref{sec:kernel}, and \(v_2>0\) by the next
display. Later division by \(w_3\) is likewise made only after
\eqref{eq:gj-Lambda} proves \(w_3>0\).
Strict log-concavity gives
\[
 v_j=\frac{u_j}{u_1}A(p_j,p_1)>0\qquad(j\geq2).
\]
The case \(k=2\) of Lemma~\ref{lem:quotient-integral} below therefore proves
strict order two directly.
At the fixed orders used in this paper, successive column division and
differentiation give the exact three-stage cascade
\begin{align}
 W(u_1,u_2,u_3)&=u_1^3W(v_2,v_3),\label{eq:W3-quotient}\\
 W(u_1,u_2,u_3,u_4)&=u_1^4W(v_2,v_3,v_4),\label{eq:W4-first}\\
 W(v_2,v_3,v_4)&=v_2^3W(w_3,w_4),\label{eq:W4-second}\\
 W(w_3,w_4)&=w_3^2\frac{\dd}{\dd t}
 \left(\frac{w_4}{w_3}\right).\label{eq:W4-third}
\end{align}
Consequently
\begin{equation}\label{eq:W4-quotient}
 W(u_1,u_2,u_3,u_4)=u_1^4v_2^3w_3^2
 \frac{\dd}{\dd t}\left(\frac{w_4}{w_3}\right).
\end{equation}
These algebraic identities require precisely the stated nonvanishing needed
to form the quotients.

Let \(\T\) denote simultaneous translation of all point arguments. From
\eqref{eq:AM},
\begin{equation}\label{eq:TlogA}
 \T A(a,b)=q(b)-q(a),\qquad \T\log A(a,b)=M(a,b).
\end{equation}
The sign convention in \eqref{eq:TlogA} will be used below. Since
\[
 \frac{v_j}{v_2}=\frac{u_j}{u_2}
 \frac{A(p_j,p_1)}{A(p_2,p_1)},
\]
direct differentiation gives
\begin{equation}\label{eq:gj-expanded}
 g_j:=\T\log(v_j/v_2)
 =A(p_j,p_2)+M(p_j,p_1)-M(p_2,p_1).
\end{equation}
The exact weighted-mean split
\[
 A(p_j,p_1)M(p_j,p_1)
 =A(p_j,p_2)M(p_j,p_2)+A(p_2,p_1)M(p_2,p_1)
\]
turns \eqref{eq:gj-expanded} into
\begin{equation}\label{eq:gj-Lambda}
 g_j=\frac{A(p_j,p_2)}{A(p_j,p_1)}
 \Lambda(p_j;p_2,p_1)>0.
\end{equation}
Thus \(w_j=(v_j/v_2)g_j>0\), and
\eqref{eq:W3-quotient} has the strict order-three sign.

Define
\[
 \Psi(\xi;m,r)=\Lambda(\xi;m,r)+\T\log\Lambda(\xi;m,r).
\]
Because \(w_j=(v_j/v_2)g_j\),
\begin{equation}\label{eq:L4-first}
 \T\log(w_4/w_3)
 =(g_4+\T\log g_4)-(g_3+\T\log g_3).
\end{equation}
Equation \eqref{eq:gj-Lambda} and \eqref{eq:TlogA} give
\[
 \T\log g_j=M(p_j,p_2)-M(p_j,p_1)+\T\log\Lambda_j,
\]
where \(\Lambda_j=\Lambda(p_j;p_2,p_1)\). A second use of the same mean split
gives
\[
 g_j+M(p_j,p_2)-M(p_j,p_1)=\Lambda_j-A(p_2,p_1).
\]
The common last term cancels in \eqref{eq:L4-first}, yielding the exact
identity
\begin{equation}\label{eq:L4-Psi}
 \T\log(w_4/w_3)
 =\Psi(p_4;p_2,p_1)-\Psi(p_3;p_2,p_1).
\end{equation}

\begin{lemma}[Iterated quotient integral]\label{lem:quotient-integral}
Let \(f_1,\ldots,f_k\) be continuously differentiable on an interval,
let \(f_1>0\), and put \(G_j=f_j/f_1\). For
\(t_1<\cdots<t_k\),
\begin{equation}\label{eq:iterated-determinant}
 \det[f_j(t_i)]_{i,j=1}^k
 =\prod_{i=1}^k f_1(t_i)
 \int_{t_1}^{t_2}\!\cdots\!\int_{t_{k-1}}^{t_k}
 \det[G'_{j+1}(\tau_i)]_{i,j=1}^{k-1}
 \dd\tau_1\cdots\dd\tau_{k-1}.
\end{equation}
The identity may be applied repeatedly whenever the first member of the new
derivative family is positive.
\end{lemma}

\begin{proof}
Factor \(f_1(t_i)\) from row \(i\). The first column is then one. Replace
successive rows, from bottom to top, by their difference with the preceding
row. Each remaining entry is
\(G_j(t_{i+1})-G_j(t_i)=\int_{t_i}^{t_{i+1}}G_j'(\tau)\dd\tau\).
Determinant multilinearity gives \eqref{eq:iterated-determinant}. Since
\(\tau_i\in[t_i,t_{i+1}]\), the integration nodes remain ordered, so the same
argument applies to the derivative determinant.
\end{proof}

\begin{lemma}[Confluent divided-difference limit]\label{lem:confluent-limit}
Let \(f\in C^{2k-2}\), let
\(V(a)=\prod_{0\leq i<j<k}(a_j-a_i)\), and let \(a,b\) be fixed strictly
increasing \(k\)-tuples. Then
\begin{equation}\label{eq:two-sided-confluence}
 \lim_{h\downarrow0}
 \frac{\det[f(z+h(a_i-b_j))]_{i,j=0}^{k-1}}
 {h^{k(k-1)}V(a)V(b)}
 =\frac{(-1)^{k(k-1)/2}}
 {\prod_{r=0}^{k-1}(r!)^2}
 \det[f^{(i+j)}(z)]_{i,j=0}^{k-1}.
\end{equation}
If only the row nodes coalesce while \(y_0<\cdots<y_{k-1}\) remain fixed,
then
\begin{equation}\label{eq:one-sided-confluence}
 \lim_{h\downarrow0}
 \frac{\det[f(z+ha_i-y_j)]_{i,j=0}^{k-1}}
 {h^{k(k-1)/2}V(a)}
 =\frac{\det[f^{(i)}(z-y_j)]_{i,j=0}^{k-1}}
 {\prod_{r=0}^{k-1}r!}.
\end{equation}
\end{lemma}

\begin{proof}
Apply successive Newton divided-difference row operations, whose divisors
are the positive numbers \(h(a_j-a_i)\), and let \(h\downarrow0\). This gives
\eqref{eq:one-sided-confluence}. Applying the same operation to the columns
gives derivatives with respect to \(-y\), hence the factor
\((-1)^{0+\cdots+(k-1)}\), and proves
\eqref{eq:two-sided-confluence}. This records the factorial normalization and
does not assume a nonzero leading coefficient.
\end{proof}

\begin{proposition}[Three-point equivalence]\label{prop:equivalence}
Under the strict lower-order positivity already proved, global \(\PF_4\) is
equivalent to
\[
 \partial_\xi\Psi(\xi;m,r)\leq0\qquad(\xi<m<r).
\]
Strict inequality implies strict global \(\PF_4\).
\end{proposition}

\begin{proof}
If \(\partial_\xi\Psi\leq0\), then \(p_4<p_3\) makes the right side of
\eqref{eq:L4-Psi} nonnegative. All prefactors in
\eqref{eq:W4-quotient} are positive. The Wronskian is therefore
nonnegative, and is positive in the strict case.

For clarity, applying Lemma~\ref{lem:quotient-integral} three times to the
order-four minor gives, with \(I_i=[t_i,t_{i+1}]\),
\begin{align*}
\det[u_j(t_i)]_{i,j=1}^4
={}&\prod_{i=1}^4u_1(t_i)
\int_{I_1\times I_2\times I_3}\prod_{i=1}^3v_2(\tau_i)\\
&\quad\times\int_{\tau_1}^{\tau_2}\int_{\tau_2}^{\tau_3}
w_3(\sigma_1)w_3(\sigma_2)
\int_{\sigma_1}^{\sigma_2}
\left(\frac{w_4}{w_3}\right)'(\rho)\dd\rho
\dd\sigma_2\dd\sigma_1\dd\boldsymbol\tau.
\end{align*}
Every factor is positive in the strict case, and every integration interval
has positive length. This proves positivity of every ordered collocation minor,
not only of the confluent Wronskian.

Conversely, assume global \(\PF_4\). Fix arbitrary
\(\xi_1<\xi_2<m<r\), choose any \(t\), and choose the columns so that
\((p_4,p_3,p_2,p_1)=(\xi_1,\xi_2,m,r)\) at \(t\); explicitly,
\(y_j=t-p_j\), which preserves \(y_1<\cdots<y_4\). Coalesce arbitrary
strictly ordered row nodes at \(t\). Equation
\eqref{eq:one-sided-confluence} and global \(\PF_4\) make the resulting
translate Wronskian nonnegative. Equations \eqref{eq:W4-quotient} and
\eqref{eq:L4-Psi} therefore give
\[
 \Psi(\xi_1;m,r)\geq\Psi(\xi_2;m,r).
\]
Because both the pair \(\xi_1<\xi_2\) and the pair \(m<r\) were arbitrary,
\(\Psi(\cdot;m,r)\) is nonincreasing for every admissible \(m,r\). Smoothness
makes this equivalent to the stated differential inequality.
\end{proof}
\section{Curvature coordinate and sign bridge}\label{sec:coordinate}

Since \(q>0\),
\[
 y(t)=-s(t),\qquad Q(y(t))=q(t)
\]
defines a global strictly increasing coordinate on its image, with
\(\dd y/\dd t=q(t)=Q(y)>0\). Primes in this and the next three sections mean
\(y\)-derivatives. Put
\[
 \rho=1-Q''=\frac{F_2}{q^3}>0,
 \qquad \kappa=2-Q''=1+\rho>1.
\]

For \(\xi<m<r\), write
\[
 p=y(\xi),\quad z=y(m),\quad w=y(r),\quad L=z-p,\quad R=w-z.
\]
Because \(A(a,b)=y(b)-y(a)\) and
\(M(a,b)=Q[y(a),y(b)]\),
\begin{equation}\label{eq:Lambda-coordinate}
 \Lambda=(w-p)+Q[p,z]-Q[z,w].
\end{equation}
Two integrations of \(\kappa=2-Q''\) give
\begin{align}
 \Lambda&=\frac1L\int_p^z(y-p)\kappa(y)\dd y
 +\frac1R\int_z^w(w-y)\kappa(y)\dd y,\label{eq:Lambda-triangular}\\
 \delta:=-\partial_p\Lambda
 &=\frac1{L^2}\int_p^z(z-y)\kappa(y)\dd y>0.
 \label{eq:delta-triangular}
\end{align}
Simultaneous translation becomes
\[
 \T=Q(p)\partial_p+Q(z)\partial_z+Q(w)\partial_w,
\]
and \(\partial_\xi=Q(p)\partial_p\). Thus
\(\Lambda_\xi=-Q(p)\delta\).

Define
\begin{equation}\label{eq:N-def}
 \mathcal N=\delta\Lambda^2+Q'(p)\delta\Lambda+
 \Lambda\T\delta-\delta\T\Lambda.
\end{equation}
Differentiating \(\Psi=\Lambda+\T\log\Lambda\), commuting simultaneous
translation with \(\partial_\xi\), and collecting over \(\Lambda^2\) gives
\begin{equation}\label{eq:sign-bridge}
 \partial_\xi\Psi=-\frac{Q(p)}{\Lambda^2}\mathcal N.
\end{equation}
Equivalently,
\begin{equation}\label{eq:K-def}
 \frac{\mathcal N}{\delta\Lambda}
 =K:=\Lambda+Q'(p)+\T\log\delta-\T\log\Lambda.
\end{equation}
It remains to prove \(K>0\).

All objects in this coordinate are evaluated on the image \(y(\R)\), where
the inverse coordinate is well defined. The displayed \(\rho,\kappa\) signs
hold at every such point, and simultaneous translation of the original
endpoints is exactly the operator \(\T\) above.
\section{The confluent invariant in curvature form}\label{sec:confluent}

Define
\begin{equation}\label{eq:D-def}
 D=3(\kappa-1)-\{Q(\log\kappa)'\}'.
\end{equation}

\begin{lemma}[Confluent-to-curvature identity]\label{lem:C4-curvature}
The determinant \eqref{eq:C4-def} satisfies
\begin{equation}\label{eq:C4-curvature}
 \Cfour=Q^6\kappa^2D.
\end{equation}
Consequently \(D>0\) on \(\R\).
\end{lemma}

\begin{proof}
The change of variable gives \(\dd/\dd t=Q\dd/\dd y\) and \(c_2=-Q\).
Recursively applying this operator produces the five cumulants listed in
\appref{app:algebra}. Substitution into the Schur-complement form of
\eqref{eq:C4-def} factors out \(Q^6\). The remaining factor is
\[
 \kappa^2\left(3(\kappa-1)-\{Q(\log\kappa)'\}'\right),
\]
as shown by the common expanded numerator in \appref{app:algebra}. This proves
\eqref{eq:C4-curvature} without a numerical assumption. Lemma
\ref{lem:certified-input}, \(Q>0\), and \(\kappa>0\) then imply \(D>0\).
\end{proof}

The determinant definition is primary: it is the doubly confluent
fourth-order translation minor after division by the positive row and column
Vandermonde factors and by \(\Phi^4\). The explicit thirteen-term expansion is
included in \appref{app:algebra} for independent checking of the certificate
input.

\section{Exact endpoint transport identity}\label{sec:transport}

Write the three normalized triangular densities from
\eqref{eq:Lambda-triangular} and \eqref{eq:delta-triangular} as
\begin{align}
 \mu_L(y)&=\frac{(z-y)\kappa(y)}{L^2\delta}
 \quad(p\leq y\leq z),\label{eq:mu}\\
 \nu_L(y)&=\frac{(y-p)\kappa(y)}{L\Lambda}
 \quad(p\leq y\leq z),\nonumber\\
 \nu_R(y)&=\frac{(w-y)\kappa(y)}{R\Lambda}
 \quad(z\leq y\leq w).\label{eq:nu}
\end{align}
The triangular identities give
\begin{equation}\label{eq:density-masses}
 \int_p^z\mu_L=1,\qquad
 \int_p^z\nu_L+\int_z^w\nu_R=1.
\end{equation}
No measure theory is needed below; these are normalized continuous weights on
compact intervals.

Set
\[
 A_0(y)=3(y-Q'(y))-Q(y)\frac{\kappa'(y)}{\kappa(y)}.
\]
Then \(A_0'=D\).

\begin{lemma}[Endpoint transport cancellation]\label{lem:transport}
With \(K\) as in \eqref{eq:K-def}, put
\begin{equation}\label{eq:transport-endpoint}
 \mathcal E=
 \frac{-I_L/L+I_R/R}{\Lambda}
 +\frac{L\mathcal H(p)-I_L}{L^2\delta},
\end{equation}
where \(I_L=\mathcal J(z)-\mathcal J(p)\) and
\(I_R=\mathcal J(w)-\mathcal J(z)\), with \(\mathcal H,\mathcal J\) defined
below. Then
\begin{equation}\label{eq:transport-expectation}
 K=\mathcal E.
\end{equation}
\end{lemma}

\begin{proof}
Two elementary primitives carry the complete cancellation:
\begin{align}
 \kappa A_0&=\mathcal H',&
 \mathcal H&=3y^2-3Q-3yQ'+(Q')^2+QQ'',\label{eq:H}\\
 \mathcal H&=\mathcal J',& \mathcal J&=y^3-3yQ+QQ'.\label{eq:J}
\end{align}
Both identities follow by one differentiation. Integrating the three linear
weights once by parts gives
\begin{equation}\label{eq:expectation-endpoint}
 \int_p^z A_0\nu_L+\int_z^w A_0\nu_R-\int_p^z A_0\mu_L
 =\mathcal E.
\end{equation}

On the other hand, differentiation of the triangular integrals along \(\T\)
gives
\[
 \T\Lambda=U(z,w)-U(p,z),
\quad
 \T\delta=\frac{U(p,z)-(Q(z)-Q(p))\delta-Q(p)\kappa(p)}L,
\]
where
\begin{equation}\label{eq:U-endpoint}
 U(c,d)=Q(c)+Q(d)-\frac{Q(d)Q'(d)-Q(c)Q'(c)}{d-c}
 +\frac{(Q(d)-Q(c))^2}{(d-c)^2}.
\end{equation}
Substitution in \eqref{eq:K-def} yields
\begin{equation}\label{eq:K-endpoint}
 K=\Lambda+Q'(p)-Q[p,z]
 +\frac{U(p,z)-U(z,w)}{\Lambda}
 +\frac{U(p,z)-Q(p)\kappa(p)}{L\delta}.
\end{equation}
The endpoint identity equating \eqref{eq:transport-endpoint} and
\eqref{eq:K-endpoint}, with every substitution exposed, is
\appref{app:endpoints}. It is an identity in the endpoint positions and the
values of \(Q,Q',Q''\).
\end{proof}
\section{The positive closed coordinate gap}\label{sec:crossing}

Define a deterministic piecewise function on \([p,w]\) by
\begin{equation}\label{eq:closed-gap}
 \Delta(y)=
 \begin{cases}
 \displaystyle
 \frac{-(z-y)^2-(z-y)Q'(y)-Q(y)+(z-p)^2+(z-p)Q'(p)+Q(p)}
      {L^2\delta}\\[3pt]
 \displaystyle\qquad
 -\frac{(y-p)^2-(y-p)Q'(y)+Q(y)-Q(p)}{L\Lambda},
 &p\leq y\leq z,\\[12pt]
 \displaystyle
 \frac{(w-y)^2+(w-y)Q'(y)+Q(y)-Q(w)}{R\Lambda},
 &z<y\leq w.
 \end{cases}
\end{equation}
The fundamental theorem of calculus and \(\kappa=2-Q''\) identify these
closed expressions with
\begin{equation}\label{eq:gap-integrals}
 \Delta(y)=
 \begin{cases}
 \displaystyle\int_p^y(\mu_L(t)-\nu_L(t))\dd t,&p\leq y\leq z,\\[6pt]
 \displaystyle\int_y^w\nu_R(t)\dd t,&z<y\leq w.
 \end{cases}
\end{equation}
The two branches agree at \(z\) by \eqref{eq:density-masses}.

\begin{lemma}[Strict closed gap]\label{lem:crossing}
The function \(\Delta\) is continuous on \([p,w]\), satisfies
\(\Delta(p)=\Delta(w)=0\), and
\[
 \Delta(y)>0\qquad(p<y<w).
\]
\end{lemma}

\begin{proof}
The endpoints and continuity follow from \eqref{eq:closed-gap},
\eqref{eq:gap-integrals}, and the exact mass identities. On \((p,z)\), the
difference of the positive left densities changes sign at the explicit strict
convex combination
\[
 y_* = \frac{\Lambda z+L\delta p}{\Lambda+L\delta}\in(p,z).
\]
Indeed, clearing the positive common factor \(\kappa(y)\) and the positive
denominators gives
\[
 \mu_L(y)-\nu_L(y)>0\iff y<y_*,
 \qquad
 \mu_L(y)-\nu_L(y)<0\iff y>y_*.
\]
Thus \(\Delta(y)>0\) for \(p<y\leq y_*\). For \(y_*\leq y\leq z\), subtract
the two identities in \eqref{eq:density-masses} and split at \(y\) to obtain
\[
 \Delta(y)=\int_y^z(\nu_L-\mu_L)\dd t+\int_z^w\nu_R\dd t>0.
\]
The first integral is nonnegative and the second is strictly positive because
\(z<w\), \(\kappa>0\), and \(\Lambda>0\). On \((z,w)\), positivity follows
directly from the positive tail integral in \eqref{eq:gap-integrals}.
\end{proof}

On the left branch \(\Delta'=\mu_L-\nu_L\), and on the right branch
\(\Delta'=-\nu_R\). Direct integration by parts on \([p,z]\) and \([z,w]\),
using \(A_0'=D\), the zero endpoint values, and the equality of the two
branches at \(z\), gives
\begin{align}
 \int_p^w\Delta(y)D(y)\dd y
 &=\int_p^zA_0\nu_L+\int_z^wA_0\nu_R-\int_p^zA_0\mu_L\nonumber\\
 &=\mathcal E=K.\label{eq:crossing-integral}
\end{align}
Combining \eqref{eq:K-def}, \eqref{eq:C4-curvature}, and
\eqref{eq:crossing-integral} gives the central identity
\begin{equation}\label{eq:central-identity}
 \mathcal N=\delta\Lambda\int_p^w
 \Delta(y)\frac{\Cfour(y)}{Q(y)^6\kappa(y)^2}\dd y>0.
\end{equation}
Every factor in the interior is strictly positive.
\section{Completion and exact order}\label{sec:completion}

Equations \eqref{eq:sign-bridge} and \eqref{eq:central-identity} imply
\[
 \partial_\xi\Psi(\xi;m,r)<0\qquad(\xi<m<r).
\]
Proposition~\ref{prop:equivalence} therefore proves part~(1) of
Theorem~\ref{thm:main}.

We now prove directly that the classification stops at order four. Let
\(h=211/2000\). Exact rational theta enclosures give
\begin{equation}\label{eq:finite-pf5-witness}
 -1.08\mathbin{\cdot}10^{-7}
 <\det[\Phi((i-j)h)]_{i,j=0}^4
 <-1.05\mathbin{\cdot}10^{-7}.
\end{equation}

\begin{theorem}\label{thm:not-pf5}
The Riemann kernel is not \(\PF_5\).
\end{theorem}

\begin{proof}
The nodes \(0,h,2h,3h,4h\) are strictly increasing. Evenness identifies the
translation matrix with the symmetric Toeplitz matrix determined by the five
kernel values \(\Phi(0),\Phi(h),\ldots,\Phi(4h)\). Each value is enclosed by
exact Taylor bounds for the retained theta modes and the infinite-tail theorem;
a rational determinant parity factorization then gives
\eqref{eq:finite-pf5-witness}. This is a negative translation minor at
distinct nodes.
\end{proof}

Equation~\eqref{eq:finite-pf5-witness} proves part~(2) of
Theorem~\ref{thm:main}. Together with the strict minors through order four,
this completes the exact-order classification.

\begin{remark}[Riemann Hypothesis boundary]
Schoenberg's transform theorem says that the bilateral Laplace transform of a
\(\PF_\infty\) function is the reciprocal of a Laguerre--P\'olya entire function
on its strip of convergence~\cite{schoenberg-ii,belton}. The Fourier transform
of \(\Phi\), up to the standard normalization, is the Riemann \(\Xi\)-function,
which has known real zeros; hence \(\Phi\) is unconditionally not
\(\PF_\infty\). The Schoenberg formulation related to the Riemann Hypothesis
instead asks for the inverse transform of the reciprocal of \(\Xi\) to be a
P\'olya-frequency function~\cite{grochenig}. Thus the finite-order
classification proved here is a different statement.
\end{remark}
\section{Verification}\label{sec:reproducibility}

The verification boundary is the theorem
\path{PF4.globalRiemannKernel_pfOrderExactly_four} in the project
\path{proof/formal}. It has two immediate mathematical inputs:
\path{PF4.globalRiemannKernel_strictPFUpTo_four}, which quantifies over every
order \(1\leq k\leq4\) and every pair of strictly increasing real node tuples,
and \path{PF4.globalRiemannKernel_orderFive_translationMinor_neg}, which proves
the negative minor at \(h=211/2000\). The former depends on the analytic theta
construction, the global \(q,F_2,\Cfour\) signs, and the quotient--transport
argument developed in Sections~\ref{sec:kernel}--\ref{sec:crossing}.

The following criteria define the formal verification claim and make it
reproducible.
\begin{enumerate}
\item \textbf{Pinned environment.} Lean \texttt{v4.32.0} is fixed in
\path{lean-toolchain}; mathlib \texttt{v4.32.0} is fixed in
\path{lake-manifest.json}. The resolved mathlib commit is
\begin{center}
{\footnotesize\texttt{81a5d257c8e410db227a6665ed08f64fea08e997}}.
\end{center}

\item \textbf{Complete build.} From \path{proof/formal}, the command
\begin{center}
\texttt{lake build}
\end{center}
elaborates every imported declaration and checks every generated proof term
with the Lean kernel. The recorded release build completes all 3737 jobs.

\item \textbf{Axiom-closure audit.} The command
\begin{center}
\texttt{lake env lean PF4/Audit.lean}
\end{center}
runs \texttt{\#print axioms} on the final theorem, both immediate inputs, and
the target-facing declarations in their proof chain. For every audited target
the printed dependency set is exactly
\[
 \{\texttt{propext},\ \texttt{Classical.choice},\ \texttt{Quot.sound}\}.
\]
These are Lean's standard logical axioms. A declaration containing
\texttt{sorry} or \texttt{admit} elaborates with the axiom
\texttt{sorryAx}; consequently such a gap anywhere in an audited target's
dependency closure would appear in this output. The observed output contains
no \texttt{sorryAx} and no project-declared axiom.

\item \textbf{Exact sign decisions.} The compact sign certificates, theta-tail
bounds, and the order-five determinant enclosure enter the theorem as exact
integer or rational inequalities. Floating-point values do not discharge a
proposition in the final dependency closure.

\item \textbf{Quantifier coverage.} The strict-order conclusion is a single
universal theorem over arbitrary real node tuples; it is not a finite node
sample. The negative order-five conclusion includes the strict ordering of
the five rational nodes, the translation-matrix identity, and the determinant
sign.
\end{enumerate}

Two independent Python replays expose the finite arithmetic printed in the
paper. From the repository root,
\begin{center}
\begin{tabular}{@{}ll@{}}
\toprule
Claim & Command \\
\midrule
Global \(q,F_2,\Cfour\) signs &
\texttt{python3 scripts/verify\_riemann\_signs\_exact.py} \\
Negative minor at \(h=211/2000\) &
\texttt{python3 scripts/verify\_pf5\_witness\_exact.py} \\
\bottomrule
\end{tabular}
\end{center}
Both use outward rational bounds. They are independently inspectable
reconstructions of the finite certificates; the universal theorem remains the
kernel-checked Lean declaration above. Tool versions, output digests, the PDF
digest, and the source-archive digest are recorded in
\path{paper/RELEASE_MANIFEST.md}.
\section{Declarations and availability}

\paragraph{Author contribution.}
Joshuah Rainstar conceived and directed the project, developed the proof
architecture, verified the exact calculations, and prepared the manuscript.

\paragraph{Acknowledgments.}
The author acknowledges OpenAI Codex and ChatGPT, Anthropic Claude, and
Moonshot AI Kimi for assistance with draft preparation and review. The author
takes responsibility for the mathematical claims and the final manuscript.

\paragraph{Conflict of interest.}
The author declares no competing interests.

\paragraph{Code and data availability.}
The public repository is
\url{https://github.com/falseywinchnet/zeta}. The submitted manuscript is
preserved by the immutable tag \texttt{pf4-paper-v1.0.0}, commit
{\footnotesize\texttt{4a14988093d7afaf05f453f8daf5441c09ba58b3}}. The present revised proof is
prepared as release \texttt{pf4-paper-v2.1.0}.
Their distinct status and material proof changes are listed in
\path{paper/REVISION_HISTORY.md}; current source, environment, verification,
PDF, and archive digests are listed in \path{paper/RELEASE_MANIFEST.md}. The
repository archive is assigned DOI
\begin{center}
\url{https://doi.org/10.5281/zenodo.21487371}.
\end{center}
\clearpage
\begin{thebibliography}{99}
\bibitem{schoenberg-ii}
I. J. Schoenberg,
\emph{On P\'olya frequency functions. II. Variation-diminishing integral
operators of the convolution type},
Acta Sci. Math. (Szeged) \textbf{12} (1950), 97--106.

\bibitem{karlin}
S. Karlin,
\emph{Total Positivity, Vol. I},
Stanford University Press, Stanford, 1968.

\bibitem{belton}
A. Belton, D. Guillot, A. Khare, and M. Putinar,
\emph{Preservers of totally positive kernels and P\'olya frequency functions},
Math. Res. Rep. \textbf{3} (2022), 35--56,
arXiv:2110.08206.

\bibitem{khare}
A. Khare,
\emph{Multiply positive functions, critical exponent phenomena, and the
Jain--Karlin--Schoenberg kernel},
arXiv:2008.05121.

\bibitem{debruijn}
N. G. de Bruijn,
\emph{The roots of trigonometric integrals},
Duke Math. J. \textbf{17} (1950), 197--226,
doi:10.1215/S0012-7094-50-01720-0.

\bibitem{dimitrov-xu}
D. K. Dimitrov and Y. Xu,
\emph{Wronskians of Fourier and Laplace transforms},
Constr. Approx. \textbf{48} (2018), 55--99,
arXiv:1606.05011.

\bibitem{csordas-varga}
G. Csordas and R. S. Varga,
\emph{Moment inequalities and the Riemann Hypothesis},
Constr. Approx. \textbf{4} (1988), 175--198,
doi:10.1007/BF02075457.

\bibitem{grochenig}
K. Grochenig,
\emph{Schoenberg's theory of totally positive functions and the Riemann zeta
function},
arXiv:2007.12889.

\bibitem{michalowski}
W. Micha\l{}owski,
\emph{On the P\'olya frequency order of the de Bruijn--Newman kernel:
certified failure at order five and the Toeplitz threshold phenomenon},
arXiv:2602.20313.
\end{thebibliography}

\appendix
\section{Confluent and curvature algebra}\label{app:algebra}

The central-moment Hankel matrix is
\[
 \begin{pmatrix}
 1&0&m_2&m_3\\
 0&m_2&m_3&m_4\\
 m_2&m_3&m_4&m_5\\
 m_3&m_4&m_5&m_6
 \end{pmatrix}.
\]
Taking the Schur complement of the upper-left entry reduces its determinant to
\[
 \det\begin{pmatrix}
 m_2&m_3&m_4\\
 m_3&m_4-m_2^2&m_5-m_2m_3\\
 m_4&m_5-m_2m_3&m_6-m_3^2
 \end{pmatrix}.
\]
Substitution of the moments from \secref{sec:inputs} and ordinary
three-by-three expansion gives
\begin{align}\label{eq:C4-expanded}
\Cfour={}&12c_2^6+24c_2^4c_4-24c_2^3c_3^2+2c_2^3c_6
-12c_2^2c_3c_5+7c_2^2c_4^2\\
&+12c_2c_3^2c_4+c_2c_4c_6-c_2c_5^2-9c_3^4-c_3^2c_6
+2c_3c_4c_5-c_4^3.
\end{align}

In the curvature coordinate, repeated application of
\(Q\dd/\dd y\) to \(c_2=-Q\) gives
\begin{align*}
c_2={}&-Q,\\
c_3={}&-QQ',\\
c_4={}&-Q\{(Q')^2+QQ''\},\\
c_5={}&-Q\{(Q')^3+4QQ'Q''+Q^2Q'''\},\\
c_6={}&-Q\{(Q')^4+11Q(Q')^2Q''+4Q^2(Q'')^2
 +7Q^2Q'Q'''+Q^3Q''''\}.
\end{align*}
Substitution in the preceding polynomial gives
\[
 \Cfour=-Q^6\mathcal P,
\]
where
\begin{align*}
\mathcal P={}&QQ''Q''''-Q(Q''')^2-2QQ''''+Q'Q''Q'''-2Q'Q'''\\
&+3(Q'')^3-15(Q'')^2+24Q''-12.
\end{align*}
Now \(\kappa=2-Q''\), and direct differentiation gives
\[
 \{Q(\log\kappa)'\}'
 =\frac{Q'\kappa'+Q\kappa''}{\kappa}
  -\frac{Q(\kappa')^2}{\kappa^2},
\]
so, after multiplication by \(\kappa^2\) and substitution of
\(\kappa'= -Q'''\), \(\kappa''=-Q''''\),
\[
 3(\kappa-1)-\{Q(\log\kappa)'\}'
 =-\frac{\mathcal P}{(Q''-2)^2}
 =-\frac{\mathcal P}{\kappa^2}.
\]
Therefore \(\Cfour=Q^6\kappa^2D\), proving
\eqref{eq:C4-curvature} by an exposed common numerator.
\section{Exact proof of the certified input}\label{app:tail}

This appendix proves Lemma~\ref{lem:certified-input}. The positive-side
theta series has \(j\)-th derivative
\[
 2\pi n^2e^{5u/2-n^2x}R_j(n^2x),\qquad x=\pi e^{2u},
\]
where the polynomials \(R_j\) satisfy the recurrence in
\secref{sec:inputs}. Machin's identity and alternating series give
\(\pi>157/50\); hence \(R_0(x)=2x-3>0\).

\subsection{Cleared sign polynomials}

In addition to \(a_j\), define the normalized remaining modes
\[
 e_j(x)=\sum_{n\geq3}
 \frac{n^2e^{-(n^2-1)x}R_j(n^2x)}{R_0(x)},
 \qquad \widehat a_j=a_j+e_j.
\]
For an indeterminate raw jet \((\widehat a_0,\ldots,\widehat a_6)\), put
\[
 N_q=\widehat a_1^2-\widehat a_0\widehat a_2,
\]
\begin{align*}
N_{F_2}={}&-\widehat a_0^4\widehat a_2\widehat a_4
+\widehat a_0^4\widehat a_3^2
+\widehat a_0^3\widehat a_1^2\widehat a_4
-2\widehat a_0^3\widehat a_1\widehat a_2\widehat a_3
+2\widehat a_0^3\widehat a_2^3\\
&-3\widehat a_0^2\widehat a_1^2\widehat a_2^2
+3\widehat a_0\widehat a_1^4\widehat a_2-\widehat a_1^6,
\end{align*}
and
\[
 N_{\Cfour}=\det[\widehat a_{i+j}]_{i,j=0}^3.
\]
Direct logarithmic differentiation and the central-moment determinant give
\begin{equation}\label{eq:raw-signs}
 q=\frac{N_q}{\widehat a_0^2},\qquad
 F_2=\frac{N_{F_2}}{\widehat a_0^6},\qquad
 \Cfour=\frac{N_{\Cfour}}{\widehat a_0^4}.
\end{equation}
The thirteen-term cumulant expansion obtained from the last identity agrees
with \eqref{eq:C4-expanded}; this is an exact polynomial identity.

For orientation, retaining only the first mode gives
\[
q_1=\frac{4x(4x^2-12x+15)}{(2x-3)^2},
\]
\[
(F_2)_1=\frac{64x^3(64x^6-576x^5+2448x^4-6432x^3
+10044x^2-8100x+2295)}{(2x-3)^6},
\]
\[
(\Cfour)_1=\frac{49152x^6(16x^4-96x^3+360x^2-840x+945)}{(2x-3)^4}.
\]
After \(x=5+z\), every numerator has positive coefficients. The second
mode is retained exactly because it supplies the larger margin near the
origin.

\subsection{Exact two-mode margin}

Rational Taylor bounds for the exponential imply
\[
 0<\eta<\frac1{12000}\quad(x\geq\pi),
\]
and
\[
 \frac1{23000}<\eta<\frac1{12000}
 \quad\left(\frac{157}{50}\leq x\leq\frac{10}{3}\right).
\]
The close endpoints are entirely rational checks. Namely,
\[
 \sum_{k=0}^{16}\frac{(471/50)^k}{k!}>12000,\qquad
 \sum_{k=0}^{21}\frac{10^k}{k!}>22000,
\]
and the geometric bound for the tail after degree \(13\) gives
\[
 e^{10}<
 \sum_{k=0}^{13}\frac{10^k}{k!}
 \frac{10^{14}/14!}{1-10/15}
 =\frac{3480060751}{154791}<23000.
\]
These prove, respectively, the upper \(\eta\)-cap, the cap
\(e^{-10}<1/22000\), and the lower bound \(e^{-10}>1/23000\).
Likewise \(\sum_{k=0}^{21}15^k/k!>3{,}000{,}000\), which gives
\(\eta<1/3{,}000{,}000\) for \(x\geq5\).
We use the following elementary certificate lemma. If
\(p(s,t)=\sum p_{ij}s^it^j\) has bidegree \((d,e)\), its Bernstein
coefficients on the unit square are
\[
 \beta_{k\ell}=\sum_{i\leq k,\,j\leq\ell}
 p_{ij}\frac{\binom{k}{i}}{\binom{d}{i}}
       \frac{\binom{\ell}{j}}{\binom{e}{j}}.
\]
The power-to-Bernstein identity proves that \(\beta_{k\ell}>0\) for every
\(k,\ell\) implies \(p>0\) on the square.

Apply this lemma after affine rational changes of variables. For \(q\), one
rectangle ends at \(x=10/3\), followed by a positive base polynomial minus
exponentially decreasing corrections. For \(F_2\), dependent bounds on the
first rectangle are followed by a rectangle through \(x=5\) and a positive
half-strip expansion. For \(\Cfour\), the first rectangle is followed by two
negative odd-power corrections whose ratios to the positive base decrease.
The derivative of each ratio has a numerator with positive coefficients after
the indicated shift. The 361 resulting rational coefficient inequalities give
\begin{equation}\label{eq:two-mode-margin}
 q_{1,2}>10,\qquad (F_2)_{1,2}>1000,\qquad
 (\Cfour)_{1,2}>50{,}000{,}000.
\end{equation}
This is a finite algebraic proof on fixed domains.

Here is the precise connection between those coefficient checks and the
quantitative margins. Evaluate the raw polynomials at the two-mode jet and set
\[
 \mathcal Q=N_q-10a_0^2,
 \qquad \mathcal F=N_{F_2}-1000a_0^6,
 \qquad \mathcal C=N_{\Cfour}-50{,}000{,}000a_0^4.
\]
The rows labelled \(Q,F,C\) below certify the cleared numerators of
\(\mathcal Q,\mathcal F,\mathcal C\), respectively. Thus their positivity is
exactly, rather than merely qualitatively, the three inequalities in
\eqref{eq:two-mode-margin}. For a ratio row, write the shifted target as
\(B_0(x)+\sum_{j\geq1}\eta^jB_j(x)\), where \(B_0\) has positive shifted
coefficients. Whenever \(B_j\) is adverse, put \(D_j=-B_j\). Positivity of
\[
 3jD_jB_0-D_j'B_0+D_jB_0'
\]
after the same shift proves that
\(e^{-3jx}D_j(x)/B_0(x)\) decreases. The displayed endpoint ratio is the sum
of all such adverse ratios at the left endpoint; being less than one proves
the target positive on the entire half-line. Coefficientwise nonnegative
pieces only increase it.

For reference, the complete finite certificate is summarized below.  Write
\(\mathcal A=(2x-3)+4(8x-3)\eta\), which is positive on every listed domain.
Rows \(Q_0,F_0,F_1,C_0\) are Bernstein certificates after the indicated
affine change to the unit square; the remaining rows use coefficient
positivity after a half-line shift. In the two ``ratio'' rows, the listed
fraction is the sum of
the adverse coefficient ratios and is strictly less than one.
\begin{center}
\tiny
\begin{tabular}{@{}lllll@{}}
\toprule
ID & Variables and domain & Degree & Positive denominator & Exact lower datum \\
\midrule
\(Q_0\) & \(157/50\le x\le10/3,\ 0\le \eta\le1/12000\)
 & \((4,2)\) & \(\mathcal A^2\) & \(6613625187767/70312500000\) \\
\(Q_\infty\) & \(x=10/3+z,\ 0\le \eta\le1/12000\)
 & \((4,2)\) & \(\mathcal A^2\) & ratio \(47333/505500\) \\
\(F_0\) & \(157/50\le x\le10/3,\ 1/23000\le \eta\le1/12000\)
 & \((12,6)\) & \(\mathcal A^6\) & minimum \(>1610047\) \\
\(F_1\) & \(10/3\le x\le5,\ 0\le \eta\le1/22000\)
 & \((12,6)\) & \(\mathcal A^6\) & \(29124973000/19683\) \\
\(F_\infty\) & \(x=5+z,\ \eta=v/3000000,\ 0\leq v\leq1\)
 & \((12,6)\) & \(\mathcal A^6\) & \(432/19073486328125\) \\
\(C_0\) & \(157/50\le x\le10/3,\ 0\le \eta\le1/12000\)
 & \((14,4)\) & \(\mathcal A^4\) & minimum \(>9546672947\) \\
\(C_\infty\) & \(x=10/3+z,\ 0\le \eta\le1/12000\)
 & \((14,4)\) & \(\mathcal A^4\) & ratio \(65989337396066791/77165720550000000\) \\
\bottomrule
\end{tabular}
\end{center}
These seven checks comprise 361 exact rational inequalities.  The derivative
envelopes used below contribute 140 further Bernstein coefficients, of
bidegrees \((j+1,1)\), \(0\le j\le6\); their smallest coefficient is
\(44861/31250\).  Thus the compact domains and the differentiated tail
majorants belong to the same exact certificate.

\subsection{All later modes}

For direct reconstruction of the derivative bounds, put
\(S_j(t)=2^j(-1)^jR_j(t)\). The recurrence gives
\begin{align*}
S_0={}&2t-3,\\
S_1={}&8t^2-30t+15,\\
S_2={}&32t^3-224t^2+330t-75,\\
S_3={}&128t^4-1440t^3+4232t^2-3270t+375,\\
S_4={}&512t^5-8448t^4+41408t^3-68096t^2+30930t-1875,\\
S_5={}&2048t^6-46592t^5+343040t^4-976320t^3
       +1008968t^2-285870t+9375,\\
S_6={}&8192t^7-245760t^6+2536960t^5-11109120t^4
       +20633312t^3-14260064t^2+2610330t-46875.
\end{align*}

For \(t\geq27\), shifting \(t=27+z\) gives positive coefficients in
\((-1)^jR_j(t)\) and in
\(2^{j+1}t^{j+1}-(-1)^jR_j(t)\). Therefore
\begin{equation}\label{eq:R-envelope}
 0<(-1)^jR_j(t)\leq2^{j+1}t^{j+1},
 \qquad 0\leq j\leq6.
\end{equation}
The derivative numerator of the normalized \(n=3\) term is another positive
shifted polynomial, so its absolute value decreases for \(x\geq3\). For
\(n\geq4\), successive majorants have ratio at most
\[
 \left(\frac54\right)^{16}e^{-27}<10^{-9}.
\]
Rational Taylor lower bounds for \(e^{24},e^{27},e^{45}\) then yield, on
\(\pi\leq x\leq5\),
\begin{equation}\label{eq:core-tail}
(|e_0|,\ldots,|e_6|)<
(7\!\cdot\!10^{-9},4\!\cdot\!10^{-7},2\!\cdot\!10^{-4},
7\!\cdot\!10^{-4},3\!\cdot\!10^{-2},1.1,41).
\end{equation}
Indeed, the normalized \(n=3\) term decreases on \(x\geq3\), and
\(e^{24}>25{,}000{,}000{,}000\), so it is bounded by
\(3|R_j(27)|/25{,}000{,}000{,}000\). The sum of the \(n\geq4\) majorants is
at most
\[
 \frac{2^{j+2}3^j4^{2j+4}}{3\cdot10^{19}}.
\]
The data entering these two expressions are
\begin{center}
\small
\begin{tabular}{@{}ccl@{}}
\toprule
\(j\) & \( |R_j(27)| \) & chosen upper bound for \( |e_j| \)\\
\midrule
0 & \(51\) & \(7\cdot10^{-9}\)\\
1 & \(5037/2\) & \(4\cdot10^{-7}\)\\
2 & \(475395/4\) & \(2\cdot10^{-4}\)\\
3 & \(42678141/8\) & \(7\cdot10^{-4}\)\\
4 & \(3623251731/16\) & \(3\cdot10^{-2}\)\\
5 & \(288709329165/32\) & \(11/10\)\\
6 & \(21373315648611/64\) & \(41\)\\
\bottomrule
\end{tabular}
\end{center}
For each row the later-mode expression is less than one thousandth of the
\(n=3\) expression, and their sum is strictly below the chosen bound. This
derives every constant in \eqref{eq:core-tail} from the printed recurrence and
three rational exponential inequalities.
The same Bernstein lemma gives
\[
(|a_0|,\ldots,|a_6|)<(2,5,20,200,2000,60000,1000000).
\]
For any one of the printed raw polynomials
\(N=\sum_\alpha c_\alpha X^\alpha\), the perturbation estimate used here is
the explicit finite sum
\begin{equation}\label{eq:perturbation-formula}
 |N(a+e)-N(a)|
 \leq\sum_\alpha |c_\alpha|
 \left\{\prod_{j=0}^6(A_j+E_j)^{\alpha_j}
              -\prod_{j=0}^6A_j^{\alpha_j}\right\},
\end{equation}
whenever \(|a_j|\leq A_j\) and \(|e_j|\leq E_j\). Substitution of the two
displayed seven-tuples into \eqref{eq:perturbation-formula} gives
\begin{equation}\label{eq:core-perturbation}
 |\Delta N_q|<0.000405,\quad
 |\Delta N_{F_2}|<37.959,\quad
 |\Delta N_{\Cfour}|<31{,}549{,}532.
\end{equation}

For \(x\geq5\), \eqref{eq:R-envelope} gives
\[
 |a_j|\leq2^{j+2}x^j,\qquad
 |e_j|\leq2^{j+2}3^{2j+4}x^je^{-8x}.
\]
The raw weights of \(N_q,N_{F_2},N_{\Cfour}\) are respectively \(2,6,12\).
Thus every perturbation majorant is a positive multiple of
\(x^we^{-8x}\), or a faster-decaying term, with \(w\leq12\); it decreases
for \(x\geq5\). In \eqref{eq:perturbation-formula} take
\[
 A_j=2^{j+2}5^j,\qquad
 E_j=\frac{2^{j+1}3^{2j+4}5^j}{10^{17}}.
\]
The latter follows from the displayed tail bound and
\(\sum_{k=0}^{47}40^k/k!>2\cdot10^{17}\). Direct substitution gives
\begin{equation}\label{eq:outer-perturbation}
 |\Delta N_q|<6.49\cdot10^{-11},\quad
 |\Delta N_{F_2}|<0.03,\quad
 |\Delta N_{\Cfour}|<418{,}287.
\end{equation}
To compare like quantities, note that
\[
 a_0=1+4\eta\frac{R_0(4x)}{R_0(x)}>1,
\]
and, by homogeneity,
\[
 N_{q,1,2}=q_{1,2}a_0^2,\qquad
 N_{F_2,1,2}=(F_2)_{1,2}a_0^6,\qquad
 N_{\Cfour,1,2}=(\Cfour)_{1,2}a_0^4.
\]
Consequently \eqref{eq:two-mode-margin} supplies the same respective lower
margins \(10\), \(1000\), and \(50{,}000{,}000\) for the cleared
numerators.  Both perturbation bounds
\eqref{eq:core-perturbation} and \eqref{eq:outer-perturbation} are strictly
smaller. Their scale relative to the corresponding cleared-numerator margin
is recorded here:
\begin{center}
\small
\begin{tabular}{@{}lcc@{}}
\toprule
target & surviving fraction on \(\pi\leq x\leq5\)
       & surviving fraction on \(x\geq5\)\\
\midrule
\(q\) & \(>0.9999595\) & \(>0.99999999999351\)\\
\(F_2\) & \(>0.962041\) & \(>0.99997\)\\
\(\Cfour\) & \(>0.36900936\) & \(>0.99163426\)\\
\bottomrule
\end{tabular}
\end{center}
Thus even the least relative margin, the compact \(\Cfour\) bound, retains
more than \(36.9\) percent of its certified two-mode numerator margin.
Finally \(\widehat a_0=a_0+e_0>0\), so \eqref{eq:raw-signs} proves all three
signs for \(u\geq0\). Evenness proves Lemma~\ref{lem:certified-input}.
\section{Endpoint algebra for the transport cancellation}\label{app:endpoints}

This appendix records the algebra suppressed by the phrase ``simplifying the
endpoint expression.'' Put
\[
 z=p+L,\qquad w=z+R,\qquad
 a=\frac{Q(z)-Q(p)}L,\qquad b=\frac{Q(w)-Q(z)}R.
\]
Direct integration of \(2-Q''\) gives
\begin{equation}\label{eq:endpoint-delta-Lambda}
 \delta=\frac{L+Q'(p)-a}{L},\qquad
 \Lambda=L+R+a-b.
\end{equation}
Define
\[
 \mathcal J_y=y^3-3yQ(y)+Q(y)Q'(y),
\]
\[
 \mathcal H_p=3p^2-3Q(p)-3pQ'(p)+Q'(p)^2+Q(p)Q''(p),
\]
and, for endpoint data,
\[
 U(c,d)=Q(c)+Q(d)-\frac{Q(d)Q'(d)-Q(c)Q'(c)}{d-c}
 +\frac{(Q(d)-Q(c))^2}{(d-c)^2}.
\]
With \(I_L=\mathcal J_z-\mathcal J_p\) and
\(I_R=\mathcal J_w-\mathcal J_z\), the identity needed in
Lemma~\ref{lem:transport} is exactly
\begin{align*}
&\frac{-(\mathcal J_z-\mathcal J_p)/L
 +(\mathcal J_w-\mathcal J_z)/R}{\Lambda}
 +\frac{L\mathcal H_p-(\mathcal J_z-\mathcal J_p)}{L^2\delta}\\
&\quad=\Lambda+Q'(p)-a
 +\frac{U(p,z)-U(z,w)}{\Lambda}
 +\frac{U(p,z)-Q(p)(2-Q''(p))}{L\delta}.
\end{align*}
To check it, substitute \eqref{eq:endpoint-delta-Lambda}, replace
\(z\) by \(p+L\) and \(w\) by \(p+L+R\), and multiply by the common
denominator \(L^2R\delta\Lambda\). The terms containing \(Q'(w)\), then
\(Q'(z)\), then \(Q''(p)\) cancel pairwise; the remaining endpoint values and
position terms cancel after inserting \(a=(Q(z)-Q(p))/L\) and
\(b=(Q(w)-Q(z))/R\). No differential equation for \(Q\) is used, so the
identity holds for arbitrary smooth \(Q\).
Equivalently, regard the endpoint jets
\(Q(p),Q(z),Q(w),Q'(p),Q'(z),Q'(w),Q''(p)\) and \(p,L,R\) as independent
indeterminates subject only to the displayed definitions of \(a,b,\delta\),
and \(\Lambda\). Clearing \(L^2R\delta\Lambda\) first and then expanding gives
the zero polynomial. Expanding the uncleared rational expression need not
expose this cancellation.
\end{document}
\typeout{get arXiv to do 4 passes: Label(s) may have changed. Rerun}
